\documentclass[10pt,a4paper]{article}
\usepackage{fontspec}
\setmainfont{Arial}
\setmonofont{Consolas}[Scale=0.92]
\usepackage[spanish,es-noshorthands]{babel}
\usepackage[a4paper,margin=1.8cm]{geometry}
\usepackage{amsmath}
\usepackage{booktabs,longtable,array}
\usepackage{xcolor}
\usepackage[hidelinks]{hyperref}
\definecolor{gris}{gray}{0.40}
\setlength{\parskip}{0.5em}
\setlength{\parindent}{0pt}
\setlength{\LTpre}{0.3em}
\title{Scripts del controlador, los simuladores, el entrenamiento y la conexión con los modelos}
\date{}
\author{}
\begin{document}
\maketitle
\vspace{-14mm}

Repositorio \texttt{wilson-cowan-id}, commit \texttt{c10e568} de la rama main. Los números de línea corresponden a esa versión.

La columna \emph{líneas} da el rango en el archivo, del encabezado (o decorador) a la última línea. La columna \emph{código} cuenta solo las líneas con código, sin comentarios, docstrings ni líneas en blanco. En gris, los métodos de cada clase.

\section*{Controlador}
El controlador del paper y el integrador del lazo.
\begin{longtable}{@{}p{0.30\textwidth}rrp{0.42\textwidth}@{}}
\toprule símbolo & líneas & código & qué hace \\ \midrule \endhead
\multicolumn{4}{@{}l}{\rule{0pt}{3.2ex}\textbf{\texttt{src/neural\_ode/closed\_loop.py}}\hfill{\small 354 líneas, 162 de código}} \\*[0.3ex]
\texttt{theta\_gamma\_refs} & 48--59 & 7 & referencia theta-gamma \\
\texttt{\_inv\_sat} & 65--73 & 4 & inversa de la sigmoidea, acotada a su dominio \\
clase \texttt{IMCController} & 76--220 & 68 & PI, inversa de la sigmoidea y desacople (port del .m) \\
\quad{\color{gris}\small \texttt{\_\_init\_\_}} & {\color{gris}\small 77--151} & {\color{gris}\small 30} & \\
\quad{\color{gris}\small \texttt{compute}} & {\color{gris}\small 153--220} & {\color{gris}\small 37} & \\
\texttt{simulate\_closed\_loop} & 304--354 & 36 & integrador original del lazo \\
\addlinespace[0.4ex]
\multicolumn{4}{@{}l}{\rule{0pt}{3.2ex}\textbf{\texttt{src/neural\_ode/closed\_loop\_ext.py}}\hfill{\small 336 líneas, 231 de código}} \\*[0.3ex]
\texttt{simulate\_closed\_loop\_ext} & 206--298 & 72 & integrador que usan todos los lazos actuales \\
\texttt{theta\_gamma\_ventana\_refs} & 304--336 & 22 & referencia theta2s (gamma con ventana theta) \\
\addlinespace[0.4ex]
\multicolumn{4}{@{}l}{\rule{0pt}{3.2ex}\textbf{\texttt{scripts/exp\_f6\_closed\_loop.py}}\hfill{\small 177 líneas, 109 de código}} \\*[0.3ex]
\texttt{fixed\_de} & 53--60 & 7 & parámetros fijos del controlador \\
\texttt{rmse} & 74--79 & 5 & error de seguimiento contra la referencia \\
\addlinespace[0.4ex]
\end{longtable}
\section*{Simuladores}
Las plantas que generan los datos y que cierran el lazo de referencia.
\begin{longtable}{@{}p{0.30\textwidth}rrp{0.42\textwidth}@{}}
\toprule símbolo & líneas & código & qué hace \\ \midrule \endhead
\multicolumn{4}{@{}l}{\rule{0pt}{3.2ex}\textbf{\texttt{src/wilson\_cowan/model.py}}\hfill{\small 600 líneas, 301 de código}} \\*[0.3ex]
clase \texttt{WilsonCowanParams} & 47--79 & 18 & $\beta$ verdaderos \\
\quad{\color{gris}\small \texttt{ke}} & {\color{gris}\small 69--74} & {\color{gris}\small 3} & \\
\quad{\color{gris}\small \texttt{ki}} & {\color{gris}\small 76--79} & {\color{gris}\small 3} & \\
\texttt{sigmoid} & 86--87 & 2 & sigmoidea \\
clase \texttt{WilsonCowan} & 302--514 & 123 & sistema de Wilson-Cowan (traducción del .m) \\
\quad{\color{gris}\small \texttt{\_\_init\_\_}} & {\color{gris}\small 311--321} & {\color{gris}\small 11} & \\
\quad{\color{gris}\small \texttt{rhs}} & {\color{gris}\small 329--343} & {\color{gris}\small 8} & \\
\quad{\color{gris}\small \texttt{rhs\_aug}} & {\color{gris}\small 356--360} & {\color{gris}\small 5} & \\
\quad{\color{gris}\small \texttt{perturbed\_field}} & {\color{gris}\small 365--395} & {\color{gris}\small 18} & \\
\quad{\color{gris}\small \texttt{effective\_inputs}} & {\color{gris}\small 399--403} & {\color{gris}\small 5} & \\
\quad{\color{gris}\small \texttt{simulate}} & {\color{gris}\small 410--459} & {\color{gris}\small 35} & \\
\quad{\color{gris}\small \texttt{\_simulate\_perturbed}} & {\color{gris}\small 464--514} & {\color{gris}\small 40} & \\
\addlinespace[0.4ex]
\multicolumn{4}{@{}l}{\rule{0pt}{3.2ex}\textbf{\texttt{src/wilson\_cowan/uncertainty.py}}\hfill{\small 634 líneas, 278 de código}} \\*[0.3ex]
clase \texttt{Perturbation} & 67--120 & 28 & clase base de las perturbaciones \\
\quad{\color{gris}\small \texttt{extra0}} & {\color{gris}\small 75--77} & {\color{gris}\small 2} & \\
\quad{\color{gris}\small \texttt{d\_extra}} & {\color{gris}\small 79--81} & {\color{gris}\small 2} & \\
\quad{\color{gris}\small \texttt{inputs}} & {\color{gris}\small 84--86} & {\color{gris}\small 2} & \\
\quad{\color{gris}\small \texttt{drive}} & {\color{gris}\small 88--90} & {\color{gris}\small 2} & \\
\quad{\color{gris}\small \texttt{gains}} & {\color{gris}\small 92--94} & {\color{gris}\small 2} & \\
\quad{\color{gris}\small \texttt{weights}} & {\color{gris}\small 96--98} & {\color{gris}\small 2} & \\
\quad{\color{gris}\small \texttt{deriv}} & {\color{gris}\small 100--102} & {\color{gris}\small 2} & \\
\quad{\color{gris}\small \texttt{sigmoid\_i}} & {\color{gris}\small 105--106} & {\color{gris}\small 2} & \\
\quad{\color{gris}\small \texttt{sigmoid\_e}} & {\color{gris}\small 108--109} & {\color{gris}\small 2} & \\
\quad{\color{gris}\small \texttt{rest\_offsets}} & {\color{gris}\small 111--113} & {\color{gris}\small 2} & \\
\quad{\color{gris}\small \texttt{metadata}} & {\color{gris}\small 115--117} & {\color{gris}\small 2} & \\
\quad{\color{gris}\small \texttt{\_\_repr\_\_}} & {\color{gris}\small 119--120} & {\color{gris}\small 2} & \\
clase \texttt{NoPerturbation} & 123--127 & 2 & planta ideal \\
clase \texttt{Refractoriness} & 137--166 & 9 & refractaria, $r = 0{,}10$ \\
\quad{\color{gris}\small \texttt{\_\_init\_\_}} & {\color{gris}\small 158--160} & {\color{gris}\small 3} & \\
\quad{\color{gris}\small \texttt{gains}} & {\color{gris}\small 162--163} & {\color{gris}\small 2} & \\
\quad{\color{gris}\small \texttt{metadata}} & {\color{gris}\small 165--166} & {\color{gris}\small 2} & \\
clase \texttt{Actuator} & 169--228 & 26 & filtro de 1 ms; con recorte=True, recorte duro en $\pm 3$ \\
\quad{\color{gris}\small \texttt{\_\_init\_\_}} & {\color{gris}\small 193--204} & {\color{gris}\small 6} & \\
\quad{\color{gris}\small \texttt{\_tau}} & {\color{gris}\small 206--209} & {\color{gris}\small 4} & \\
\quad{\color{gris}\small \texttt{d\_extra}} & {\color{gris}\small 211--214} & {\color{gris}\small 3} & \\
\quad{\color{gris}\small \texttt{inputs}} & {\color{gris}\small 216--222} & {\color{gris}\small 5} & \\
\quad{\color{gris}\small \texttt{metadata}} & {\color{gris}\small 224--228} & {\color{gris}\small 5} & \\
\addlinespace[0.4ex]
\multicolumn{4}{@{}l}{\rule{0pt}{3.2ex}\textbf{\texttt{src/neural\_ode/closed\_loop.py}}\hfill{\small 354 líneas, 162 de código}} \\*[0.3ex]
\texttt{make\_perturbed\_plant} & 255--275 & 9 & simulador perturbado como planta del lazo \\
\addlinespace[0.4ex]
\multicolumn{4}{@{}l}{\rule{0pt}{3.2ex}\textbf{\texttt{scripts/gen\_uncertain\_dataset.py}}\hfill{\small 189 líneas, 102 de código}} \\*[0.3ex]
\texttt{generar\_con} & 106--173 & 53 & simula y guarda un dataset \\
\addlinespace[0.4ex]
\multicolumn{4}{@{}l}{\rule{0pt}{3.2ex}\textbf{\texttt{scripts/esc\_gen\_amplitud.py}}\hfill{\small 198 líneas, 124 de código}} \\*[0.3ex]
\texttt{hz} & 58--59 & 2 & frecuencias de cada familia \\
\texttt{ventana} & 62--64 & 2 & ventana de encendido \\
\texttt{square} & 69--70 & 2 & familia square \\
\texttt{aprbs} & 73--85 & 12 & familia APRBS \\
\texttt{thetagamma} & 88--94 & 6 & familia theta-gamma \\
\texttt{poisson} & 97--112 & 14 & familia Poisson \\
\texttt{chirp} & 115--121 & 6 & familia chirp \\
\texttt{ruido} & 124--136 & 12 & ruido blanco de banda limitada \\
\texttt{par} & 141--162 & 20 & parámetros de cada escenario según familia y amplitud \\
\texttt{escenarios} & 165--181 & 16 & grilla de amplitudes 1 a 12 y particiones train, val y test \\
\texttt{main} & 184--194 & 11 & genera los datasets de las tres plantas \\
\addlinespace[0.4ex]
\end{longtable}
\section*{Entrenamiento}
De los datos al checkpoint. Los modelos en sí (\texttt{GrayBoxWC} y \texttt{LagGrayBox}) están en la sección siguiente, porque son los mismos que entran al lazo.
\begin{longtable}{@{}p{0.30\textwidth}rrp{0.42\textwidth}@{}}
\toprule símbolo & líneas & código & qué hace \\ \midrule \endhead
\multicolumn{4}{@{}l}{\rule{0pt}{3.2ex}\textbf{\texttt{scripts/esc\_gen\_ruido.py}}\hfill{\small 85 líneas, 47 de código}} \\*[0.3ex]
\texttt{agregar\_ruido} & 39--63 & 23 & suma ruido de observación al dataset ya generado (σ 0,01 y 0,05, por semilla) \\
\addlinespace[0.4ex]
\multicolumn{4}{@{}l}{\rule{0pt}{3.2ex}\textbf{\texttt{scripts/cola\_amplitud.py}}\hfill{\small 95 líneas, 59 de código}} \\*[0.3ex]
\texttt{args\_de} & 57--64 & 7 & argumentos de cada configuración \\
\texttt{tareas} & 67--90 & 24 & las 180 corridas e25, e26 y e27, con su punto de partida \\
\addlinespace[0.4ex]
\multicolumn{4}{@{}l}{\rule{0pt}{3.2ex}\textbf{\texttt{scripts/cola.py}}\hfill{\small 96 líneas, 60 de código}} \\*[0.3ex]
\texttt{correr} & 50--96 & 36 & cola de corridas con lugares fijos; lanza esc\_run.py por tarea \\
\addlinespace[0.4ex]
\multicolumn{4}{@{}l}{\rule{0pt}{3.2ex}\textbf{\texttt{scripts/esc\_run.py}}\hfill{\small 273 líneas, 198 de código}} \\*[0.3ex]
\texttt{cargar\_params} & 49--59 & 8 & warm start: carga $\hat\beta$ de otra corrida (opción init-params) \\
\texttt{suavizar} & 62--77 & 11 & suavizado opcional de los datos (no se usa en las corridas actuales) \\
\texttt{main} & 80--269 & 158 & una corrida: arma el modelo, elige fit o fit\_aug, evalúa y guarda .pt y .json \\
\addlinespace[0.4ex]
\multicolumn{4}{@{}l}{\rule{0pt}{3.2ex}\textbf{\texttt{src/neural\_ode/graybox\_train.py}}\hfill{\small 548 líneas, 321 de código}} \\*[0.3ex]
\texttt{apilar\_historia} & 46--56 & 6 & ventana de comandos pasados para las historias \\
\texttt{make\_windows} & 59--98 & 27 & corta las trayectorias en ventanas de entrenamiento \\
\texttt{open\_loop\_mse} & 167--187 & 18 & error en corrida libre, para validación \\
\texttt{param\_errors} & 190--194 & 4 & error de $\hat\beta$ contra los verdaderos \\
clase \texttt{MejorVal} & 197--226 & 21 & guarda el mejor modelo por validación (incluye el de inicio) \\
\quad{\color{gris}\small \texttt{\_\_init\_\_}} & {\color{gris}\small 204--209} & {\color{gris}\small 6} & \\
\quad{\color{gris}\small \texttt{evaluar}} & {\color{gris}\small 211--219} & {\color{gris}\small 8} & \\
\quad{\color{gris}\small \texttt{restaurar}} & {\color{gris}\small 221--226} & {\color{gris}\small 6} & \\
clase \texttt{TrainConfig} & 233--257 & 25 & hiperparámetros de fit \\
\texttt{build\_model} & 291--300 & 8 & construye el GrayBoxWC según la variante \\
\texttt{fit} & 303--506 & 142 & entrenamiento de white-box, agnóstica, historias, filtro entrenable y black-box \\
\texttt{load\_split} & 518--548 & 24 & lee el dataset con sus particiones \\
\addlinespace[0.4ex]
\multicolumn{4}{@{}l}{\rule{0pt}{3.2ex}\textbf{\texttt{src/neural\_ode/memory.py}}\hfill{\small 310 líneas, 186 de código}} \\*[0.3ex]
clase \texttt{AugTrainConfig} & 198--211 & 14 & hiperparámetros de fit\_aug, incluido x0 entrenable \\
\texttt{fit\_aug} & 214--310 & 73 & entrenamiento del backbone expandido del actuador (estado aumentado) \\
\addlinespace[0.4ex]
\multicolumn{4}{@{}l}{\rule{0pt}{3.2ex}\textbf{\texttt{src/neural\_ode/integrate.py}}\hfill{\small 43 líneas, 16 de código}} \\*[0.3ex]
\texttt{rk4\_step} & 21--27 & 6 & paso RK4 en torch \\
\texttt{rollout} & 30--43 & 8 & integra el modelo sobre una ventana; el gradiente lo atraviesa \\
\addlinespace[0.4ex]
\multicolumn{4}{@{}l}{\rule{0pt}{3.2ex}\textbf{\texttt{scripts/esc\_eval.py}}\hfill{\small 296 líneas, 194 de código}} \\*[0.3ex]
\texttt{nrmse\_test} & 151--172 & 20 & NRMSE en corrida libre sobre test \\
\texttt{r2\_delta} & 200--244 & 38 & R² de la corrección contra el término faltante \\
\texttt{evaluar} & 250--281 & 30 & las métricas que se guardan en el .json \\
\addlinespace[0.4ex]
\end{longtable}
\section*{Conexión entre el controlador y los modelos}
Cómo cada modelo entrenado reemplaza al simulador bajo el controlador.
\begin{longtable}{@{}p{0.30\textwidth}rrp{0.42\textwidth}@{}}
\toprule símbolo & líneas & código & qué hace \\ \midrule \endhead
\multicolumn{4}{@{}l}{\rule{0pt}{3.2ex}\textbf{\texttt{scripts/control\_paper\_planta\_modelo.py}}\hfill{\small 277 líneas, 208 de código}} \\*[0.3ex]
clase \texttt{CampoModelo} & 41--100 & 53 & campo del modelo en numpy ($\hat\beta$ + red, o solo red) \\
\quad{\color{gris}\small \texttt{\_\_init\_\_}} & {\color{gris}\small 44--65} & {\color{gris}\small 21} & \\
\quad{\color{gris}\small \texttt{\_\_call\_\_}} & {\color{gris}\small 67--84} & {\color{gris}\small 17} & \\
\quad{\color{gris}\small \texttt{\_red}} & {\color{gris}\small 86--100} & {\color{gris}\small 14} & \\
clase \texttt{PlantaHistoria} & 103--134 & 19 & guarda el comando pasado para los modelos con historia \\
\quad{\color{gris}\small \texttt{\_\_init\_\_}} & {\color{gris}\small 107--110} & {\color{gris}\small 4} & \\
\quad{\color{gris}\small \texttt{extra0}} & {\color{gris}\small 112--113} & {\color{gris}\small 2} & \\
\quad{\color{gris}\small \texttt{\_\_call\_\_}} & {\color{gris}\small 115--134} & {\color{gris}\small 12} & \\
clase \texttt{PlantaFiltro} & 137--157 & 16 & estado $\hat P$ y recorte del backbone expandido \\
\quad{\color{gris}\small \texttt{\_\_init\_\_}} & {\color{gris}\small 141--145} & {\color{gris}\small 5} & \\
\quad{\color{gris}\small \texttt{extra0}} & {\color{gris}\small 147--148} & {\color{gris}\small 2} & \\
\quad{\color{gris}\small \texttt{\_\_call\_\_}} & {\color{gris}\small 150--157} & {\color{gris}\small 8} & \\
\texttt{planta\_de} & 160--166 & 7 & elige la clase de planta según el modelo \\
\texttt{verificar} & 169--193 & 24 & compara el campo numpy contra PyTorch \\
\addlinespace[0.4ex]
\multicolumn{4}{@{}l}{\rule{0pt}{3.2ex}\textbf{\texttt{scripts/esc\_eval.py}}\hfill{\small 296 líneas, 194 de código}} \\*[0.3ex]
\texttt{cargar} & 49--77 & 27 & lee el checkpoint y reconstruye el modelo \\
\addlinespace[0.4ex]
\multicolumn{4}{@{}l}{\rule{0pt}{3.2ex}\textbf{\texttt{src/neural\_ode/dynamics.py}}\hfill{\small 351 líneas, 182 de código}} \\*[0.3ex]
clase \texttt{\_FIRMLP} & 29--46 & 8 & red del filtro entrenable \\
\quad{\color{gris}\small \texttt{\_\_init\_\_}} & {\color{gris}\small 39--42} & {\color{gris}\small 4} & \\
\quad{\color{gris}\small \texttt{forward}} & {\color{gris}\small 44--46} & {\color{gris}\small 3} & \\
clase \texttt{GrayBoxWC} & 49--351 & 170 & white-box, agnóstica, historias, filtro entrenable, black-box \\
\quad{\color{gris}\small \texttt{\_\_init\_\_}} & {\color{gris}\small 55--184} & {\color{gris}\small 61} & \\
\quad{\color{gris}\small \texttt{weights}} & {\color{gris}\small 187--188} & {\color{gris}\small 2} & \\
\quad{\color{gris}\small \texttt{weights\_dict}} & {\color{gris}\small 190--192} & {\color{gris}\small 3} & \\
\quad{\color{gris}\small \texttt{\_extra}} & {\color{gris}\small 195--196} & {\color{gris}\small 2} & \\
\quad{\color{gris}\small \texttt{params\_dict}} & {\color{gris}\small 198--204} & {\color{gris}\small 6} & \\
\quad{\color{gris}\small \texttt{backbone}} & {\color{gris}\small 211--257} & {\color{gris}\small 30} & \\
\quad{\color{gris}\small \texttt{structured\_dict}} & {\color{gris}\small 260--266} & {\color{gris}\small 7} & \\
\quad{\color{gris}\small \texttt{g\_out}} & {\color{gris}\small 271--289} & {\color{gris}\small 18} & \\
\quad{\color{gris}\small \texttt{forward}} & {\color{gris}\small 295--301} & {\color{gris}\small 7} & \\
\quad{\color{gris}\small \texttt{backbone\_sensitivities}} & {\color{gris}\small 317--351} & {\color{gris}\small 30} & \\
\addlinespace[0.4ex]
\multicolumn{4}{@{}l}{\rule{0pt}{3.2ex}\textbf{\texttt{src/neural\_ode/memory.py}}\hfill{\small 310 líneas, 186 de código}} \\*[0.3ex]
clase \texttt{LagGrayBox} & 46--135 & 48 & backbone expandido del actuador ($\hat\tau$, $1/\hat\alpha$) \\
\quad{\color{gris}\small \texttt{\_\_init\_\_}} & {\color{gris}\small 60--81} & {\color{gris}\small 11} & \\
\quad{\color{gris}\small \texttt{tau}} & {\color{gris}\small 83--84} & {\color{gris}\small 2} & \\
\quad{\color{gris}\small \texttt{\_eff}} & {\color{gris}\small 86--90} & {\color{gris}\small 5} & \\
\quad{\color{gris}\small \texttt{h0}} & {\color{gris}\small 92--96} & {\color{gris}\small 2} & \\
\quad{\color{gris}\small \texttt{filtered\_inputs}} & {\color{gris}\small 98--115} & {\color{gris}\small 9} & \\
\quad{\color{gris}\small \texttt{forward}} & {\color{gris}\small 117--127} & {\color{gris}\small 11} & \\
\quad{\color{gris}\small \texttt{params\_dict}} & {\color{gris}\small 129--130} & {\color{gris}\small 2} & \\
\quad{\color{gris}\small \texttt{extras\_dict}} & {\color{gris}\small 132--135} & {\color{gris}\small 4} & \\
\addlinespace[0.4ex]
\multicolumn{4}{@{}l}{\rule{0pt}{3.2ex}\textbf{\texttt{scripts/control\_distancia\_y\_rango.py}}\hfill{\small 260 líneas, 183 de código}} \\*[0.3ex]
\texttt{modelos} & 64--98 & 31 & qué modelos entran por planta (sufijos @sat3, @satinf) \\
\texttt{ref\_en\_rango} & 117--146 & 25 & referencias en rango \\
\texttt{referencia} & 149--157 & 9 & las cuatro referencias \\
\texttt{una} & 164--203 & 38 & una corrida de lazo: simulador o modelo, controlador, integración \\
\texttt{distancias} & 210--220 & 10 & distancia entre lazos \\
\addlinespace[0.4ex]
\multicolumn{4}{@{}l}{\rule{0pt}{3.2ex}\textbf{\texttt{scripts/control\_amplitud.py}}\hfill{\small 68 líneas, 49 de código}} \\*[0.3ex]
\texttt{tareas} & 41--52 & 12 & lista de corridas de lazo (e25, e26, e27) \\
\addlinespace[0.4ex]
\multicolumn{4}{@{}l}{\rule{0pt}{3.2ex}\textbf{\texttt{scripts/trazas\_lazo\_referencias.py}}\hfill{\small 98 líneas, 80 de código}} \\*[0.3ex]
\texttt{una} & 53--83 & 30 & una corrida, con la traza completa y la referencia \\
\addlinespace[0.4ex]
\end{longtable}

\section*{Scripts viejos}
\sloppy
No intervienen en los resultados actuales.
\begin{itemize}\setlength{\itemsep}{0pt}
\item \texttt{control\_refrac.py}
\item \texttt{control\_act.py}
\item \texttt{control\_comparado.py}
\item \texttt{cola\_control.py}
\item \texttt{control\_paper\_modelos.py}
\item \texttt{control\_historia\_pasada.py}
\item \texttt{eval\_closed\_loop.py}
\end{itemize}
De \texttt{control\_paper\_modelos.py} solo se importan constantes, y de \texttt{exp\_f6\_closed\_loop.py} solo \texttt{fixed\_de} y \texttt{rmse}. La clase \texttt{IMCExtendido} de \texttt{closed\_loop\_ext.py} invierte el modelo y no se usa, porque los resultados dejan el controlador del paper sin cambios. El MATLAB original del simulador está en \texttt{documentos/simulador\_wilson\_cowan.m}.
\end{document}
